Write \(B=B_{J_y,J_z,n}\), \(K=K_{J_y,J_z,n}\), \(M=2^{dJ_z}\), and \(D=D_{J_y,J_z}\).  All constants below are uniform over the model class.  This uniformity follows from the two density bounds and the fixed wavelet regularity constants.

\emph{Step 1: localization after both Gram conjugations.}
The covariate wavelets through level \(J_z\) span the terminal boundary scaling space.  Use its \(L^2(dz)\)-orthonormal scaling basis \((\varphi_\kappa)_\kappa\), and put
\[
h_z=2^{-J_z},\qquad
G_z=(\langle\varphi_\kappa,r_n\varphi_\lambda\rangle)_{\kappa\lambda}.
\]
The covariate-density bounds give \(cI\preceq G_z\preceq CI\).  Compact support gives a fixed integer \(b_0\), independent of \(J_z\), such that \((G_z)_{\kappa\lambda}=0\) whenever the terminal-grid distance \({\rm dist}(\kappa,\lambda)>b_0\).  Set \(a=(C+c)/2\), \(R=I-G_z/a\), and \(q_0=(C-c)/(C+c)<1\).  Then
\[
G_z^{-1}=a^{-1}\sum_{j=0}^{\infty}R^j,
\qquad \|R\|_{\rm op}\le q_0,
\qquad (R^j)_{\kappa\lambda}=0\quad
\text{if }j<{\rm dist}(\kappa,\lambda)/b_0.
\]
Consequently
\[
|(G_z^{-1})_{\kappa\lambda}|
\le \frac{a^{-1}}{1-q_0}
q_0^{{\rm dist}(\kappa,\lambda)/b_0-1}
\le C_0e^{-c_0{\rm dist}(\kappa,\lambda)}.                 \tag{1}
\]
Thus this decay is obtained from the convergent Neumann series; it does not assert that the inverse Gram matrix is compactly supported.

Let
\[
L_{J_z}(z,u)=\sum_{\kappa,\lambda}
\varphi_\kappa(z)(G_z^{-1})_{\kappa\lambda}\varphi_\lambda(u).
\]
At any point only a fixed number of scaling functions are nonzero, each is bounded by \(C h_z^{-d/2}\), and terminal-grid distance dominates \(h_z^{-1}\|z-u\|_2\) up to a fixed additive constant.  Equation (1) therefore gives
\[
|L_{J_z}(z,u)|\le C h_z^{-d}e^{-c\|z-u\|_2/h_z}
=CM e^{-cM^{1/d}\|z-u\|_2},
\qquad
\sup_z\int |L_{J_z}(z,u)|du\le C.                           \tag{2}
\]

The tensor Gram matrix in the raw feature coordinates is \(I_y\otimes G_z\), because the outcome features are orthonormal for \(dy\).  Hence the two Gram inverses in \(B=G_R^{-1}H^{\rm form}G_R^{-1}\) act through two copies of (2).  For fixed \(u\), the weighted Neumann inverse has the energy identity
\[
\int f(y)\{\mathcal G_{n,u}g\}(y)dy
=\int_0^1\frac{F_f(t)F_g(t)}{\bar q_n(t,u)}dt,
\qquad F_f(t)=\int_0^t f(v)dv,                              \tag{3}
\]
for zero-mean \(f,g\).  Indeed, integrating the defining differential equation and using the Neumann boundary condition gives
\(\partial_y\mathcal G_{n,u}f=-F_f/\bar q_n\), and integration by parts gives (3).  Since \(c\le\bar q_n\le C\), the integral kernel in (3) is uniformly bounded and Lipschitz in each outcome argument.  Projection through \(J_y\) preserves a uniform bound: at outcome level \(j\), cancellation of a nonconstant wavelet against a Lipschitz function gives a coefficient bounded by \(C2^{-3j/2}\); multiplication by the pointwise wavelet bound \(C2^{j/2}\), followed by summation over \(j\), gives \(C\sum_j2^{-j}<\infty\).  Applying this in the two outcome arguments yields
\[
\sup_{u,y,y'}|G^y_{J_y,u}(y,y')|\le C.                       \tag{4}
\]

Let \(b_J(w,w')=\Phi_{J_y,J_z}(w)^\top B\Phi_{J_y,J_z}(w')\), where \(w=(y,z)\).  Equations (2)--(4), now after both Gram conjugations, give the nonseparable representation and bound
\[
b_J((y,z),(y',z'))
=\int L_{J_z}(z,u)G^y_{J_y,u}(y,y')L_{J_z}(z',u)r_n(u)du,
\]
\[
|b_J((y,z),(y',z'))|
\le CM e^{-cM^{1/d}\|z-z'\|_2}.                              \tag{5}
\]
The convolution bound in (5) uses
\(\|z-u\|_2+\|z'-u\|_2\ge\|z-z'\|_2\); a harmless polynomial factor is absorbed by reducing \(c\).  Integrating (5), and using the joint-density upper bound, proves for \(p=2,4\)
\[
\iint |b_J(w,w')|^p dP_a(w)dP_{a'}(w')\le C M^{p-1}.          \tag{6}
\]
If \(\widetilde b_{aa'}\) denotes centering of \(b_J\) in both arguments under \(P_a\) and \(P_{a'}\), conditional expectation is an \(L^p\)-contraction.  Hence (6) also holds with \(b_J\) replaced by \(\widetilde b_{aa'}\).  This completes the required \(L^2\) and \(L^4\) control after both, rather than before, Gram inversions.

\emph{Step 2: the exact edge array and its variance.}
Index an evaluation observation by \(r=(e,a,i)\), and put \(Z_r=\Phi(W_r)-E\Phi(W_r)\).  Since \(n=2m\) and the two fold values receive weight one half, direct expansion of the diagonal-free correction gives, for unordered pairs,
\[
h_{n,rs}(x,y)=
\begin{cases}
2(m-1)^{-1}x^\top By,&e_r=e_s,\ a_r=a_s,\\
-2m^{-1}x^\top By,&e_r=e_s,\ a_r\ne a_s,\\
0,&e_r\ne e_s,
\end{cases}
\qquad
n\mathbb U_{J_y,J_z,n}=\sum_{r<s}h_{n,rs}(Z_r,Z_s).          \tag{7}
\]
Independence from the two-arm sampling assumption and \(EZ_r=0\) make every kernel in (7) canonical in either argument.  Equations (6)--(7) imply
\[
\max_r\sum_{s\ne r}Eh_{n,rs}^2\le C\frac{M}{m},
\qquad
\sum_{r<s}Eh_{n,rs}^4\le C\frac{M^3}{m^2}.                   \tag{8}
\]
For one fold, canonicality kills covariances between distinct edges, so its variance is
\[
\frac{2m}{m-1}\sum_{a=0}^1
{\rm tr}(B\Gamma_aB\Gamma_a)
+4{\rm tr}(B\Gamma_0B\Gamma_1)
=2{\rm tr}(K^2)\{1+o(1)\}.                                  \tag{9}
\]
The folds are independent.  Thus (9) and the definition of \(A_{J_y,J_z,n}\) give
\[
{\rm Var}(n\mathbb U_{J_y,J_z,n})
=4{\rm tr}(K^2)\{1+o(1)\}=A_{J_y,J_z,n}^2\{1+o(1)\}.        \tag{10}
\]
The spectral proposition gives
\[
\frac{{\rm tr}(K^4)}{{\rm tr}(K^2)^2}
\le\frac{\|K\|_{\rm op}^2}{\|K\|_{\rm HS}^2}
\le CM^{-1}.                                                  \tag{11}
\]

\emph{Step 3: fourth moments and Lindeberg.}
Order the observations, let \(\mathcal F_r\) be their filtration, write
\[
R_r=\sum_{s<r}h_{n,sr}(Z_s,Z_r),
\qquad \mathcal D_{n,r}=\ell_{n,r}(W_r)+R_r.
\]
Canonicality in the new argument yields
\(E(\mathcal D_{n,r}\mid\mathcal F_{r-1})=0\), and summation gives \(\sum_r\mathcal D_{n,r}=G_n+n\mathbb U_{J_y,J_z,n}\).  Given \(Z_r\), the terms in \(R_r\) are independent and centered.  Expanding their fourth power, rather than replacing it by the desired conclusion, gives
\[
ER_r^4
=\sum_{s<r}Eh_{n,sr}^4
+6\sum_{s<t<r}E\!\left[E(h_{n,sr}^2\mid Z_r)
E(h_{n,tr}^2\mid Z_r)\right].                                \tag{12}
\]
Equation (5) gives
\(\sup_w E\{h_{n,sr}(Z_s,w)^2\}\le CM/m^2\).  Summing (12) in \(r\), using (8), gives
\[
\sum_rER_r^4\le C\left(\frac{M^3}{m^2}+\frac{M^2}{m}\right).\tag{13}
\]
The uniformly bounded linear summands satisfy
\[
\sum_rE\ell_{n,r}^4
\le C\sum_rE\ell_{n,r}^2
=CEG_n^2\le CS_n^2.                                           \tag{14}
\]
Because \(A^2\asymp M\) and \(S_n\ge A\), (13)--(14) and
\((x+y)^4\le8(x^4+y^4)\) yield
\[
S_n^{-4}\sum_rE\mathcal D_{n,r}^4
\le C\left(M^{-1}+m^{-1}+\frac{M}{m^2}\right).               \tag{15}
\]
For every \(\epsilon>0\), the elementary inequality
\(x^2\mathbf 1\{|x|>\epsilon\}\le\epsilon^{-2}x^4\), applied to \(x=\mathcal D_{n,r}/S_n\), proves the displayed Lindeberg conclusion from (15).  Here \(M\to\infty\) because \(J_z\to\infty\), and \(M\le D=o(n^2)\) gives \(M/m^2\to0\).

\emph{Step 4: predictable-variance contractions.}
Expanding the conditional square gives exactly
\[
\sum_rE(\mathcal D_{n,r}^2\mid\mathcal F_{r-1})
=\sum_rE\ell_{n,r}^2
+2\sum_r\sum_{s<r}E\{\ell_{n,r}(W_r)h_{n,sr}(Z_s,W_r)\mid Z_s\}
\]
\[
\hspace{18mm}
+\sum_r\sum_{s,t<r}E\{h_{n,sr}(Z_s,W_r)h_{n,tr}(Z_t,W_r)\mid Z_s,Z_t\}. \tag{16}
\]
The first term is deterministic.  To expand the last term, define
\[
a_{sr}(x)=E\{h_{n,sr}(x,W_r)^2\},
\qquad
c_{st}^{(r)}(x,y)=E\{h_{n,sr}(x,W_r)h_{n,tr}(y,W_r)\}.
\]
Its centered diagonal part is \(\sum_s\sum_{r>s}\{a_{sr}(Z_s)-Ea_{sr}(Z_s)\}\).  The pointwise bound after (12) gives \(\|a_{sr}\|_\infty\le CM/m^2\); grouping by \(s\), of which there are \(O(m)\), and noting that each group has \(O(m)\) terms proves
\[
E\left|\sum_s\sum_{r>s}\{a_{sr}(Z_s)-Ea_{sr}(Z_s)\}\right|^2
\le C m\left(m\frac{M}{m^2}\right)^2=C\frac{M^2}{m}.        \tag{17a}
\]
The off-diagonal part is \(2\sum_{s<t}\sum_{r>t}c_{st}^{(r)}(Z_s,Z_t)\).  Each \(c_{st}^{(r)}\) is canonical in \(x\) and in \(y\).  When its centered square is expanded, a pair of terms sharing exactly one old vertex has a remaining vertex that appears once and hence has expectation zero.  Identical or reversed old edges survive; expansion of their two contractions over the later vertices produces a four-edge cycle.  Endpoint and \((m-1)^{-1}-m^{-1}\) coefficient corrections are bounded by the fourth-moment sum in (8), namely \(CM^3/m^2\).  A genuine four-cycle with arm labels \(a_1,\ldots,a_4\) equals, after whitening,
\[
{\rm tr}\{K T_{a_1}K T_{a_2}K T_{a_3}K T_{a_4}\},
\qquad
T_a=\Gamma^{-1/2}\Gamma_a\Gamma^{-1/2},\quad0\preceq T_a\preceq I.
\]
Schatten H\"older and \(0\preceq T_a\preceq I\) bound its absolute value by \({\rm tr}(K^4)\).  The edge coefficients in (7) sum to a uniformly bounded weight for each ordered cycle, so all four-cycle contributions, including reversed orientations, are at most \(C{\rm tr}(K^4)\).  Equations (17a) and (8), the exact vanishing of the one-shared-vertex patterns, and the four-cycle calculation give
\[
E\left|\{\text{last term of (16)}\}
-E\{\text{last term of (16)}\}\right|^2
\le C\left(\frac{M^2}{m}+\frac{M^3}{m^2}+{\rm tr}(K^4)\right). \tag{17}
\]

It remains to control the middle, mixed linear--quadratic term.  Define the contraction
\[
(\mathcal T_{sr}f)(x)=E\{h_{n,sr}(x,W_r)f(W_r)\}.
\]
Whitening the bilinear form shows \(\|\mathcal T_{sr}\|_{2\to2}\le C\|K\|_{\rm op}/m\).  For a fixed old vertex, Cauchy--Schwarz in the at most \(2m\) later vertices, with the factor \(m^{-1}\) retained, bounds the squared grouped contraction by \(C\|K\|_{\rm op}^2m^{-1}\) times the sum of their linear variances.  Summing over old vertices counts each linear variance at most \(2m\) times; independence of the linear summands therefore gives
\[
E\left|2\sum_r\sum_{s<r}(\mathcal T_{sr}\ell_{n,r})(Z_s)\right|^2
\le C\|K\|_{\rm op}^2EG_n^2\le CS_n^2.                      \tag{18}
\]
Equation (18) is the mixed contraction; it is not absorbed into a pure quadratic bound.  Divide (17)--(18) by \(S_n^4\), use \(S_n^2\ge A^2\asymp M\), (11), and Cauchy--Schwarz for the covariance between the centered middle and last terms.  This yields
\[
\frac{E\left[\left\{\sum_rE(\mathcal D_{n,r}^2\mid\mathcal F_{r-1})
-E(G_n+n\mathbb U_{J_y,J_z,n})^2\right\}^2\right]}{S_n^4}
\le C\left(m^{-1}+\frac{M}{m^2}+M^{-1}\right)+o(1)\to0.      \tag{19}
\]
Equations (5)--(6) establish the requested nonseparable kernel bounds after both Gram conjugations; (12), (17), and (18) explicitly account for the fourth-moment, repeated-edge, shared-vertex, four-cycle, and mixed contractions.  Equations (15) and (19) prove the two asserted martingale conditions, completing the proof.